% TOP_PROBLEM: {"id": 30001143, "title": "Variational Principles for Dilute Spin Systems", "category_id": 19, "category": {"id": 19, "name": "probability", "display_name": "Probability", "description": "Problems involving probability theory, stochastic processes, and random structures.", "slug": "probability", "order_index": 19, "created_at": "2026-07-31T15:26:25.671Z"}, "status": "open"}
% FIRST_POSTED: null
% ENTRY_KIND: statement_only
% !TeX program = lualatex
\documentclass[11pt]{article}
\usepackage[margin=25mm]{geometry}
\usepackage{fontspec}
\setmainfont{Latin Modern Roman}
\usepackage{amsmath,amssymb,mathtools}
\usepackage{unicode-math}
\setmathfont{Latin Modern Math}
\usepackage{xurl,hyperref}
\hypersetup{colorlinks=true,urlcolor=blue}
\setcounter{secnumdepth}{0}
\setlength{\parindent}{0pt}
\setlength{\parskip}{5pt}
\setlength{\emergencystretch}{3em}
\newcommand{\R}{\mathbb R}
\newcommand{\N}{\mathbb N}
\newcommand{\Z}{\mathbb Z}
\newcommand{\Q}{\mathbb Q}
\newcommand{\C}{\mathbb C}
\newcommand{\D}{\mathbb D}
\newcommand{\F}{\mathbb F}
\newcommand{\sref}[2]{\hyperlink{source:#1:#2}{[S#2]}}
\newcommand{\cataloguescope}[1]{#1}
\begin{document}
% UnsolvedMath ID 30001143; source catalogue code OWR-3386-002
\setcounter{section}{30001142}
\section{Variational Principles for Dilute Spin Systems}\label{30001143}
{\small UnsolvedMath ID 30001143 · Probability}
\subsection{Definitions and mathematical statement}

\paragraph{Shared notation.}$\mathbb N$, $\mathbb Z$, $\mathbb Q$, $\mathbb R$, and $\mathbb C$ denote the natural numbers, integers, rationals, reals, and complex numbers, respectively. All additional parameters and hypotheses are specified in the question.


\paragraph{Statement recorded in the supplied source.}
For dilute ferromagnets and sparse spin glasses, does the GREM-associated trial order parameter give the correct variational formula for the pressure? Does the numerator have a direct differential-equation formulation analogous to the Parisi equation? For the dilute ferromagnet, can one prove $\log Z_{N+M}\leq\log Z_N+\log Z_M$ directly? How should the pressure variational principle be formulated to exhibit appropriate trial order parameters, an entropy-minus-internal-energy interpretation, or a Legendre structure?
\par\smallskip\textit{Formulation references:} \sref{30001143}{1}, \sref{30001143}{2}.
\subsection{Short English statement}
For dilute ferromagnets and sparse spin glasses, does the GREM-associated trial order parameter give the correct variational formula for the pressure? Does the numerator have a direct differential-equation formulation analogous to the Parisi equation? For the dilute ferromagnet, can one prove $\log Z_{N+M}\leq\log Z_N+\log Z_M$ directly? How should the pressure variational principle be formulated to exhibit appropriate trial order parameters, an entropy-minus-internal-energy interpretation, or a Legendre structure?
\subsection{Sources}
\begin{itemize}
\item[\hypertarget{source:30001143:1}{S1}] ULAM.AI and the UnsolvedMath Contributors, UnsolvedMath: A Curated Collection of Open Mathematics Problems, record 30001143 (OWR-3386-002). Catalogue attribution under CC BY 4.0.
\url{https://huggingface.co/datasets/ulamai/UnsolvedMath}.
\item[\hypertarget{source:30001143:2}{S2}] Original problem formulation and mathematical context.
\url{https://doi.org/10.4171/owr/2009/03}.
\end{itemize}
\label{end:30001143}
\end{document}
